% Shared preamble of the proof overview documents  <Folder>/<Name>_overview.tex.
% Each overview starts with
%     % Shared preamble of the proof overview documents  <Folder>/<Name>_overview.tex.
% Each overview starts with
%     % Shared preamble of the proof overview documents  <Folder>/<Name>_overview.tex.
% Each overview starts with
%     \input{../../scripts/overview_preamble.tex}
% (for ThinMatching/: \input{../scripts/overview_preamble.tex}) and is compiled from its own folder by
% scripts/build_overviews.py.
\documentclass[10pt,a4paper]{article}
\usepackage[T1]{fontenc}
\usepackage[utf8]{inputenc}
\usepackage[margin=2.2cm]{geometry}
\usepackage{amsmath,amssymb,amsthm}
\usepackage{array,booktabs,longtable}
\usepackage[dvipsnames]{xcolor}
\usepackage{enumitem}
\usepackage{tikz}
\usetikzlibrary{arrows.meta,positioning,shapes.multipart,fit,calc}
\usepackage[hidelinks]{hyperref}
\setlist{itemsep=2pt,topsep=4pt}
\renewcommand{\arraystretch}{1.2}

% Lean names, files and paper references
\newcommand{\lean}[1]{\textcolor{blue!50!black}{\texttt{\detokenize{#1}}}}
\newcommand{\file}[1]{\textcolor{OliveGreen!70!black}{\texttt{\detokenize{#1}}}}
\newcommand{\paperref}[1]{\textsf{#1}}
\newcommand{\hyp}[1]{\textcolor{BrickRed}{\textbf{#1}}}

% Boxes for the dependency diagrams
\tikzset{
  res/.style={draw, rounded corners, fill=blue!6, align=center, inner sep=4pt, font=\small},
  main/.style={res, fill=blue!18, very thick},
  cited/.style={res, fill=red!10, dashed},
  shared/.style={res, fill=gray!12},
  dep/.style={-{Stealth[length=2.2mm]}, thick},
}

\newcommand{\overviewtitle}[2]{%
  \begin{center}{\Large\bfseries #1}\\[4pt]{\small Proof overview of the Lean~4/Mathlib formalisation}\\[2pt]
  {\small #2}\end{center}}

% Standard paragraph describing the paper
\newcommand{\paperversion}{\noindent\textbf{Paper.} J.~Breitung, T.~Gervens, M.~Grohe.
  \emph{Graph Similarity Meets Optimal Transport}.\\ File \file{Graph_Similarity_Meets_Optimal_Transport.pdf} in the
  repository; statement, equation, page and line numbers refer to this PDF. \file{docs/NUMBERING.md} lists the Lean
  files and declarations of all numbered statements.\par\smallskip}

% (for ThinMatching/: \input{../scripts/overview_preamble.tex}) and is compiled from its own folder by
% scripts/build_overviews.py.
\documentclass[10pt,a4paper]{article}
\usepackage[T1]{fontenc}
\usepackage[utf8]{inputenc}
\usepackage[margin=2.2cm]{geometry}
\usepackage{amsmath,amssymb,amsthm}
\usepackage{array,booktabs,longtable}
\usepackage[dvipsnames]{xcolor}
\usepackage{enumitem}
\usepackage{tikz}
\usetikzlibrary{arrows.meta,positioning,shapes.multipart,fit,calc}
\usepackage[hidelinks]{hyperref}
\setlist{itemsep=2pt,topsep=4pt}
\renewcommand{\arraystretch}{1.2}

% Lean names, files and paper references
\newcommand{\lean}[1]{\textcolor{blue!50!black}{\texttt{\detokenize{#1}}}}
\newcommand{\file}[1]{\textcolor{OliveGreen!70!black}{\texttt{\detokenize{#1}}}}
\newcommand{\paperref}[1]{\textsf{#1}}
\newcommand{\hyp}[1]{\textcolor{BrickRed}{\textbf{#1}}}

% Boxes for the dependency diagrams
\tikzset{
  res/.style={draw, rounded corners, fill=blue!6, align=center, inner sep=4pt, font=\small},
  main/.style={res, fill=blue!18, very thick},
  cited/.style={res, fill=red!10, dashed},
  shared/.style={res, fill=gray!12},
  dep/.style={-{Stealth[length=2.2mm]}, thick},
}

\newcommand{\overviewtitle}[2]{%
  \begin{center}{\Large\bfseries #1}\\[4pt]{\small Proof overview of the Lean~4/Mathlib formalisation}\\[2pt]
  {\small #2}\end{center}}

% Standard paragraph describing the paper
\newcommand{\paperversion}{\noindent\textbf{Paper.} J.~Breitung, T.~Gervens, M.~Grohe.
  \emph{Graph Similarity Meets Optimal Transport}.\\ File \file{Graph_Similarity_Meets_Optimal_Transport.pdf} in the
  repository; statement, equation, page and line numbers refer to this PDF. \file{docs/NUMBERING.md} lists the Lean
  files and declarations of all numbered statements.\par\smallskip}

% (for ThinMatching/: \input{../scripts/overview_preamble.tex}) and is compiled from its own folder by
% scripts/build_overviews.py.
\documentclass[10pt,a4paper]{article}
\usepackage[T1]{fontenc}
\usepackage[utf8]{inputenc}
\usepackage[margin=2.2cm]{geometry}
\usepackage{amsmath,amssymb,amsthm}
\usepackage{array,booktabs,longtable}
\usepackage[dvipsnames]{xcolor}
\usepackage{enumitem}
\usepackage{tikz}
\usetikzlibrary{arrows.meta,positioning,shapes.multipart,fit,calc}
\usepackage[hidelinks]{hyperref}
\setlist{itemsep=2pt,topsep=4pt}
\renewcommand{\arraystretch}{1.2}

% Lean names, files and paper references
\newcommand{\lean}[1]{\textcolor{blue!50!black}{\texttt{\detokenize{#1}}}}
\newcommand{\file}[1]{\textcolor{OliveGreen!70!black}{\texttt{\detokenize{#1}}}}
\newcommand{\paperref}[1]{\textsf{#1}}
\newcommand{\hyp}[1]{\textcolor{BrickRed}{\textbf{#1}}}

% Boxes for the dependency diagrams
\tikzset{
  res/.style={draw, rounded corners, fill=blue!6, align=center, inner sep=4pt, font=\small},
  main/.style={res, fill=blue!18, very thick},
  cited/.style={res, fill=red!10, dashed},
  shared/.style={res, fill=gray!12},
  dep/.style={-{Stealth[length=2.2mm]}, thick},
}

\newcommand{\overviewtitle}[2]{%
  \begin{center}{\Large\bfseries #1}\\[4pt]{\small Proof overview of the Lean~4/Mathlib formalisation}\\[2pt]
  {\small #2}\end{center}}

% Standard paragraph describing the paper
\newcommand{\paperversion}{\noindent\textbf{Paper.} J.~Breitung, T.~Gervens, M.~Grohe.
  \emph{Graph Similarity Meets Optimal Transport}.\\ File \file{Graph_Similarity_Meets_Optimal_Transport.pdf} in the
  repository; statement, equation, page and line numbers refer to this PDF. \file{docs/NUMBERING.md} lists the Lean
  files and declarations of all numbered statements.\par\smallskip}

\begin{document}
\overviewtitle{Theorem 18: hardness of approximating the optimal-transport cut distance}{Folder \file{GsimOT/Theorem22_CutHardness}}
\paperversion

\section{Statement}

\paragraph{Paper.} Theorem~18 (Section~4, p.~13, l.~443--444): \emph{for every $\varepsilon>0$, unless
$\mathrm{P}=\mathrm{NP}$, there is no polynomial-time approximation algorithm computing $d^{\mathrm{OT}}_\square$
with multiplicative error guarantee $17/16-\varepsilon$}; this is Theorem~4\,(3) (p.~4, l.~138--139). The proof
(App.~G, pp.~57--62, l.~1826--1988) reduces from MaxCut. For a graph $F$ on $n$ vertices and $m$ edges it takes the
graphs $G',H'$ of [GG22, Theorem~23] on $n'=2m+n$ vertices with $C_\square(\pi)=\|A_{G'}-A_{H'}\|_\square=
2\,\mathrm{maxcut}(F)$ for the canonical bijection $\pi\colon v_i\mapsto w_i$, attaches $\ell_i=i(n^2+n)k$
leaves to the $i$-th vertex of both graphs, and shows, for a fixed $0<\delta<1$ and $k=\lceil1600n^6/\delta\rceil$
(p.~62, l.~1977, l.~1983),
\[ 2(1-\delta)\,\mathrm{maxcut}(F)\le d^{\mathrm{OT}}_\square(G,H)\le2\,\mathrm{maxcut}(F). \]
The case $\delta=1/34$ is Eq.~(47) (p.~58, l.~1838), $\tfrac{33}{17}\,\mathrm{maxcut}(F)\le
d^{\mathrm{OT}}_\square(G,H)\le2\,\mathrm{maxcut}(F)$. With [H\aa s01] (MaxCut cannot be approximated within
$16/17+\varepsilon$ unless $\mathrm{P}=\mathrm{NP}$), a $c$-approximation of $d^{\mathrm{OT}}_\square$ with
$c<17/16$ would approximate $\mathrm{maxcut}(F)$ within $c/(1-\delta)<17/16$ (l.~1984--1988).

\paragraph{Lean.} The correctness of the reduction is formalised.
\begin{itemize}
\item \lean{GsimOT.Hardness.maxcut_dOTCut_bounds_paper}: for every $0<\delta<1$ and every graph $F$ on
  \lean{Fin n} with an edge, with $G'=$ \lean{baseG F}, $H'=$ \lean{baseH F} and
  $\ell_i=$ \lean{leavesPaper} $F\,\delta\,i=i(n^2+n)\lceil1600n^6/\delta\rceil$,
  \[ 2(1-\delta)\,\mathrm{maxcut}(F)\le d^{\mathrm{OT}}_\square(G,H)\ \wedge\ d^{\mathrm{OT}}_\square(G,H)\le2\,\mathrm{maxcut}(F). \]
  This is the paper's construction with the paper's constants.
\item \lean{GsimOT.Hardness.maxcut_dOTCut_bounds_eq47}: Eq.~(47), the case $\delta=1/34$.
\item \lean{GsimOT.Hardness.le_dOTCut_addLeaves}, \lean{GsimOT.Hardness.dOTCut_addLeaves_le}: the lower and the
  upper bound for arbitrary base graphs on $n'\le n^2$ vertices with maximum degree $\le n-1$ and
  $\|A_{G'}-A_{H'}\|_\square=2c$, $1\le c\le n^2/4$ (the properties the proof uses).
\item \lean{GsimOT.Hardness.maxcut_dOTCut_bounds}: the variant with $n'$ in place of $n$ in the leaf counts.
\end{itemize}
$d^{\mathrm{OT}}_\square=$ \lean{dOTCut} is the infimum over $\mathrm{DS}(V,W)$ of $\|B(G,H,Q)\|_\square$
(\file{Common/CutNorm.lean}); \lean{maxcut F} is the maximum number of edges cut by a vertex set.

\section{Where the proof lives}
\begin{longtable}{>{\raggedright\arraybackslash}p{5.2cm} >{\raggedright\arraybackslash}p{10.6cm}}
\toprule \textbf{File} & \textbf{Content} \\ \midrule\endhead \bottomrule\endfoot
\file{BaseGraphs.lean} & the base graphs $G',H'$ built from the [AN06] matrix (\lean{anG}, \lean{anH}),
  \lean{maxcut}, and $\|A_{G'}-A_{H'}\|_\square=2\,\mathrm{maxcut}(F)$ (\lean{cutNorm_anG_sub_anH}). \\
\file{ProductSums.lean} & \lean{sum4}: the entry sum of $B(G,H,Q)$ over $(S_1\times S_2)\times(T_1\times T_2)$,
  a lower bound for $\|B\|_\square$ (\lean{abs_sum4_le_cutNorm}). \\
\file{Claim56_LeafMass.lean} & Claim~58 (\lean{claim56}). \\
\file{Claim57_SplitMass.lean} & Claim~59 (\lean{claim57}). \\
\file{Claim58_WrongVertex.lean} & Claim~60 (\lean{claim58}, \lean{claim58'}). \\
\file{Claim59_NearIdentity.lean} & Claim~61 (\lean{claim59}, \lean{claim59_delta}). \\
\file{Theorem22.lean} & the upper bound, \lean{leafCount}, the case distinction and the choice of constants
  (\lean{le_dOTCut_addLeaves}), the facts on $G',H',F$ (degrees, $n'\le n^2$, $\mathrm{maxcut}\le n^2/4$),
  the complete reductions and Eq.~(47). \\
\file{Common/Leaves.lean}, \file{Common/CutNorm.lean}, \file{Common/DoublyStochastic.lean} &
  \lean{addLeaves}, \lean{leavesOf}, \lean{cutNorm}, \lean{dOTCut}, \lean{mass}, \lean{le_iInf_isDS}. \\
\end{longtable}

\section{Proof layout}
\begin{center}
\begin{tikzpicture}
  \node[main] (thm) at (0,0) {Theorem 18, reduction (paper's construction)\\\lean{maxcut_dOTCut_bounds_paper}};
  \node[res] (eq47) at (5.6,0) {Eq.~(47), $\delta=1/34$\\\lean{maxcut_dOTCut_bounds_eq47}};
  \node[res] (up) at (-5.0,-2.1) {upper bound\\\lean{dOTCut_addLeaves_le}};
  \node[res] (low) at (0,-2.1) {lower bound, $k=\lceil1600n^6/\delta\rceil$\\\lean{le_dOTCut_addLeaves}};
  \node[res] (facts) at (5.2,-2.1) {$\Delta\le n-1$, $n'\le n^2$, $4\,\mathrm{maxcut}\le n^2$\\\lean{card_adj_anG_le}, \lean{card_adj_anH_le}\\\lean{two_mul_card_edge_add_le}\\\lean{four_mul_maxcut_le}};
  \node[res] (c58) at (-5.4,-4.5) {Claim 58\\\lean{claim56}};
  \node[res] (c59) at (-2.0,-4.5) {Claim 59\\\lean{claim57}};
  \node[res] (c60) at (1.6,-4.5) {Claim 60\\\lean{claim58}, \lean{claim58'}};
  \node[res] (c61) at (5.2,-4.5) {Claim 61\\\lean{claim59}, \lean{claim59_delta}};
  \node[res] (s4) at (-1.6,-6.5) {product sums\\\lean{sum4}, \lean{abs_sum4_le_cutNorm}};
  \node[res] (gg) at (4.4,-6.5) {base graphs [GG22]/[AN06]\\\lean{cutNorm_anG_sub_anH}};
  \node[shared] (com) at (-1.6,-8.4) {cut norm, $\mathrm{DS}$, leaves (Common)\\\lean{cutNorm}, \lean{mass}, \lean{addLeaves}};
  \draw[dep] (eq47) -- (thm);
  \draw[dep] (thm) -- (up); \draw[dep] (thm) -- (low); \draw[dep] (thm) -- (facts);
  \draw[dep] (low) -- (c58); \draw[dep] (low) -- (c59); \draw[dep] (low) -- (c60); \draw[dep] (low) -- (c61);
  \draw[dep] (c58) -- (s4); \draw[dep] (c59) -- (s4); \draw[dep] (c60) -- (s4); \draw[dep] (c61) -- (s4);
  \draw[dep] (facts.south east) to[out=-60,in=20] (gg.east);
  \draw[dep] (s4) -- (com);
\end{tikzpicture}
\end{center}
Blue: this folder; grey: \file{GsimOT/Common}. Nothing is cited. (\lean{maxcut_dOTCut_bounds_paper} also
uses \lean{cutNorm_anG_sub_anH} for the hypothesis $\|A_{G'}-A_{H'}\|_\square=2\,\mathrm{maxcut}(F)$.)

\section{The proof step by step}

\paragraph{The base graphs (p.~57, l.~1826--1834; p.~59, l.~1885--1891).} The proof takes $G',H'$ from
[GG22, Theorem~23]; the proof of Claim~59 describes them through the $2m\times n$ matrix of [AN06,
Proposition~3.2], in which every edge of $F$ contributes two rows with one entry $+1$ and one entry $-1$: $G'$ has
the edges of the positive entries, $H'$ those of the negative entries. \file{BaseGraphs.lean} builds exactly
these graphs: the vertices are the signed edges $e^\pm$ and the vertices of $F$; for $e=uv$, $u<v$, $G'$ has the
edges $e^+u$, $e^-v$ and $H'$ the edges $e^+v$, $e^-u$. Then $A_{G'}-A_{H'}$ consists of the [AN06] matrix and its
transpose, and \lean{cutNorm_anG_sub_anH} proves $\|A_{G'}-A_{H'}\|_\square=2\,\mathrm{maxcut}(F)$. The leaves
are \lean{addLeaves} with $\ell_i=$ \lean{leafCount n k i} $=(i+1)(n^2+n)k$ (vertices numbered from $0$).

\paragraph{Upper bound (p.~58, l.~1839--1840).} \lean{dOTCut_addLeaves_le}. The identity coupling gives
$\|B(G,H,P_{\mathrm{id}})\|_\square\le\|A_G-A_H\|_\square$ (\lean{cutNorm_B_refl_le}), and attaching the same
leaves to both graphs does not change $A_G-A_H$ on the base (\lean{cutNorm_addLeaves_sub_le}).

\paragraph{Claim~58 (pp.~58--59, l.~1854--1882).} \lean{claim56}: if $Q(v,L_H)\ge\eta$, then
$\|B\|_\square\ge\eta(M-2)/4$; the two cases ($Q(v,L(w))\le\eta/2$ for all $w$, or $>\eta/2$ for some $w^*$) and
the sets $S',T'$ are the paper's. The conclusion ($\ge\frac18\eta M\ge n^2/2\ge2\,\mathrm{maxcut}(F)$ for
$k\ge4/\eta$) is carried out with the concrete constants in \lean{le_dOTCut_addLeaves}.

\paragraph{Claim~59 (pp.~59--60, l.~1883--1905).} \lean{claim57}: if $Q(v,w),Q(v,w')\ge\tau$ for distinct base
vertices, then $\|B\|_\square\ge\tau(M-2d)/2$, where $d$ bounds the degrees in $H'$; the proof uses
$|I|\le d=n-1$ (l.~1885--1892), which is \lean{card_adj_anH_le}.

\paragraph{Reduction to near-bijections (p.~60, l.~1906--1909).} In \lean{le_dOTCut_addLeaves}: if neither claim
applies (with $\eta=\alpha/2$, $\tau=\alpha/(4n^2)$), each base vertex $v$ sends mass
$\ge1-(n'-1)\tau-\eta\ge1-\alpha$ to a single base vertex $f(v)$ (\lean{hf}).

\paragraph{Claim~60 (p.~60, l.~1913--1932).} If $f(i)\ne i$: \lean{claim58} (for $\ell_i<\ell_j$, with
$S'=\{(v_i,w_j)\}$, $T'=(V\setminus N(v_i))\times L(w_j)$, as in the paper) and \lean{claim58'} (for
$\ell_i>\ell_j$, the paper's ``reversing the roles of the two graphs'') give
$\|B\|_\square\ge(1-\rho)(|\ell_j-\ell_i|-d)$ with $d=n-1$; the gap $|\ell_j-\ell_i|\ge(n^2+n)k$ is \lean{hgap}.

\paragraph{Claim~61 (p.~61, l.~1938--1965).} If $f=\mathrm{id}$: \lean{claim59} takes the maximiser $(S,T)$ of
$\|A_{G'}-A_{H'}\|_\square$ (\lean{exists_eq_cutNorm}) and $S'=S\times W$, $T'=T\times W$, and shows
$\|B\|_\square\ge\|A_{G'}-A_{H'}\|_\square-2\alpha n'^2$, estimating the main part and the error part together for
each pair $(s,t)$. \lean{claim59_delta} derives Claim~61 as stated: for $0\le\alpha\le\delta/(2n')^2$,
$\|B\|_\square\ge2(1-\delta)\,\mathrm{maxcut}(F)$.

\paragraph{Choice of the constants (pp.~61--62, l.~1966--1977).} \lean{le_dOTCut_addLeaves} uses the paper's
$\alpha=\delta/(200n^4)$, $\eta=\alpha/2$, $\tau=\alpha/(4n^2)$, $k=\lceil1600n^6/\delta\rceil$ and $n'\le n^2$,
and checks each case: Claims~58--60 give $\|B\|_\square\ge n^2/2\ge2c$, Claim~61 gives
$\|B\|_\square\ge2c-\delta/100\ge2(1-\delta)c$ (using $c\ge1$). The infimum over $Q$ is \lean{le_iInf_isDS}.

\paragraph{The facts on $F$ (pp.~59--62).} For the construction above, \lean{card_adj_anG_le},
\lean{card_adj_anH_le} (maximum degree $\le n-1$; l.~1885--1892), \lean{two_mul_card_edge_add_le}
($n'=2m+n\le n^2$; l.~1969) and \lean{four_mul_maxcut_le} ($\mathrm{maxcut}(F)\le n^2/4$; l.~1880) are proved,
and \lean{maxcut_dOTCut_bounds_paper} combines them with \lean{le_dOTCut_addLeaves}; reindexing to
\lean{Fin n'} is \lean{baseG}, \lean{baseH} (\lean{card_adj_reindexGraph}, \lean{cutNorm_baseG_sub_baseH}).
\lean{maxcut_dOTCut_bounds_eq47} is the case $\delta=1/34$, where $2(1-\delta)=33/17$.

\section{Assumptions and trusted base}
No cited results and no hypotheses beyond $0<\delta<1$ and ``$F$ has an edge''. Only the standard axioms
\lean{propext}, \lean{Classical.choice}, \lean{Quot.sound}. The complexity side ([H\aa s01], the polynomial size
of $G,H$, l.~1977--1988) is not formalised.

\section{Notes on the formalisation}
\begin{itemize}
\item The base graphs are built in \file{BaseGraphs.lean} as described in the proof of Claim~59, with a proof of
  $\|A_{G'}-A_{H'}\|_\square=2\,\mathrm{maxcut}(F)$.
\item Claims~58--61 are stated for general parameters; the thresholds on $k$ are checked once, in
  \lean{le_dOTCut_addLeaves}, with the paper's constants.
\item \lean{maxcut_dOTCut_bounds} uses $n'$ in place of $n$ in the leaf counts and needs neither the degree bound
  nor $n'\le n^2$.
\end{itemize}

\section{How to check this result}
\begin{itemize}
\item \texttt{lake build GsimOT.Theorem22\_CutHardness.Theorem22}
\item \texttt{lake env lean scripts/CheckAxioms.lean} prints the axioms of
  \lean{GsimOT.Hardness.maxcut_dOTCut_bounds_paper} and the claims.
\item The table \file{Theorem22_correspondence.pdf} in this folder lists every paper item with its Lean
  declaration and line.
\end{itemize}
\end{document}
